\documentclass[11pt]{article}
\usepackage[a4paper,margin=18mm]{geometry}
\usepackage[no-math]{fontspec}
\setmainfont{Latin Modern Roman}
\usepackage{amsmath,amssymb,amsthm,xfrac,xspace,hyperref,graphicx,comment}
\excludecomment{DefTralics}
\providecommand{\C}{}
\providecommand{\bibliofont}{\small}
\newcommand{\ie}{i.e.,\xspace}
\newcommand{\eg}{e.g.,\xspace}
\newcommand{\etc}{etc.\xspace}
\newcommand{\cf}{cf.\xspace}
\newcommand{\resp}{resp.\xspace}
\newcommand{\smaller}{\small}
\newcommand{\Subsection}{\subsection}
\newcommand{\Subsubsection}{\subsubsection}
\newcommand{\skpt}{\smallskip}
\emergencystretch=3em
\providecommand{\loccit}{\textit{loc.\ cit.}}
\input{author-preamble.tex}
\usepackage{etoolbox}
\makeatletter
\crefname{Thm}{theorem}{theorems}
\Crefname{Thm}{Theorem}{Theorems}
\crefname{Prop}{proposition}{propositions}
\Crefname{Prop}{Proposition}{Propositions}
\crefname{Lem}{lemma}{lemmas}
\Crefname{Lem}{Lemma}{Lemmas}
\crefname{Cor}{corollary}{corollarys}
\Crefname{Cor}{Corollary}{Corollarys}
\crefname{Def}{definition}{definitions}
\Crefname{Def}{Definition}{Definitions}
\crefname{Exa}{example}{examples}
\Crefname{Exa}{Example}{Examples}
\crefname{Exas}{examples}{exampless}
\Crefname{Exas}{Examples}{Exampless}
\crefname{Cons}{construction}{constructions}
\Crefname{Cons}{Construction}{Constructions}
\crefname{Hyp}{hypothesis}{hypothesiss}
\Crefname{Hyp}{Hypothesis}{Hypothesiss}
\crefname{Rem}{remark}{remarks}
\Crefname{Rem}{Remark}{Remarks}
\crefname{Rec}{recollection}{recollections}
\Crefname{Rec}{Recollection}{Recollections}
\makeatother
\makeatletter
\expandafter\def\csname jeplabeltype@Hyp:big-tt\endcsname{Hyp}
\expandafter\def\csname jeplabeltype@Thm:support-intro\endcsname{Thm}
\expandafter\def\csname jeplabeltype@Def:weak-ring\endcsname{Def}
\expandafter\def\csname jeplabeltype@Thm:rings-intro\endcsname{Thm}
\expandafter\def\csname jeplabeltype@Thm:Spc-image-intro\endcsname{Thm}
\expandafter\def\csname jeplabeltype@Rec:MT\endcsname{Rec}
\expandafter\def\csname jeplabeltype@Rec:tensor-on-MT\endcsname{Rec}
\expandafter\def\csname jeplabeltype@Rec:A/C\endcsname{Rec}
\expandafter\def\csname jeplabeltype@Prop:A/C\endcsname{Prop}
\expandafter\def\csname jeplabeltype@Rem:Sigma\endcsname{Rem}
\expandafter\def\csname jeplabeltype@Cons:E_C\endcsname{Cons}
\expandafter\def\csname jeplabeltype@Rec:A/B\endcsname{Rec}
\expandafter\def\csname jeplabeltype@Prop:A/B\endcsname{Prop}
\expandafter\def\csname jeplabeltype@Lem:Ker[-,E]\endcsname{Lem}
\expandafter\def\csname jeplabeltype@Rec:Spch\endcsname{Rec}
\expandafter\def\csname jeplabeltype@Rem:simple\endcsname{Rem}
\expandafter\def\csname jeplabeltype@Prop:ring-in-AB\endcsname{Prop}
\expandafter\def\csname jeplabeltype@Rem:max-vs-max\endcsname{Rem}
\expandafter\def\csname jeplabeltype@Prop:EB=EC\endcsname{Prop}
\expandafter\def\csname jeplabeltype@Prop:Ker@<C<Ker[]\endcsname{Prop}
\expandafter\def\csname jeplabeltype@Thm:max-to-max\endcsname{Thm}
\expandafter\def\csname jeplabeltype@Prop:hB-hBp\endcsname{Prop}
\expandafter\def\csname jeplabeltype@Def:Supp\endcsname{Def}
\expandafter\def\csname jeplabeltype@Prop:Supp\endcsname{Prop}
\expandafter\def\csname jeplabeltype@Prop:support\endcsname{Prop}
\expandafter\def\csname jeplabeltype@Prop:agree\endcsname{Prop}
\expandafter\def\csname jeplabeltype@Thm:TPF\endcsname{Thm}
\expandafter\def\csname jeplabeltype@Thm:rings\endcsname{Thm}
\expandafter\def\csname jeplabeltype@Exa:Supp(EB)\endcsname{Exa}
\expandafter\def\csname jeplabeltype@Cor:[EB,EC]\endcsname{Cor}
\expandafter\def\csname jeplabeltype@Rec:T-S\endcsname{Rec}
\expandafter\def\csname jeplabeltype@Prop:Finv\endcsname{Prop}
\expandafter\def\csname jeplabeltype@Lem:Finv\endcsname{Lem}
\expandafter\def\csname jeplabeltype@Thm:Spc(F)\endcsname{Thm}
\expandafter\def\csname jeplabeltype@Thm:Spc-image\endcsname{Thm}
\expandafter\def\csname jeplabeltype@Thm:hom-local\endcsname{Thm}
\AtBeginDocument{\let\jeporiginalLabel\label
\renewcommand{\label}[1]{\ifcsname jeplabeltype@#1\endcsname
\edef\jeplabeltype{\csname jeplabeltype@#1\endcsname}\expandafter\jeporiginalLabel\expandafter[\jeplabeltype]{#1}\else\jeporiginalLabel{#1}\fi}}
\makeatother
\begin{document}
\begin{center}{\Large Stable item 2032}\end{center}
\paragraph{Source and attribution.}
Paul Balmer, \emph{Homological support of big objects in tensor-triangulated categories}, Journal de l’École polytechnique — Mathématiques 7 (2020), publisher version, DOI 10.5802/jep.135. \href{https://jep.centre-mersenne.org/articles/10.5802/jep.135/}{Publisher article}. Authors' text: \href{https://creativecommons.org/licenses/by/4.0/}{CC BY 4.0}.
\paragraph{Editorial problem boundary.}
Prove only Theorem 1.2 below for a rigidly-compactly generated tensor-triangulated category: the constructed homological support has the listed compact agreement, coproduct, triangle/suspension and tensor-product properties for arbitrary objects. All new injective-hull comparison and maximal-localising tensor-ideal core, the exact support construction and complete proofs of its properties are retained. No injectivity or bijectivity of the homological-to-triangular-spectrum map is assumed.
\paragraph{Difficulty (editorial): H2.}
Standard Day convolution, Gabriel localisation and pure-injective representability do not by themselves identify the maximal localising tensor ideals needed for this support. The author establishes the unique maximal localising tensor ideal extending a given maximal Serre ideal, using the internal-Hom kernel of its pure-injective residue weak ring and injective-hull comparisons across quotients. This nonroutine correspondence supplies the vanishing dichotomy that proves the tensor-product formula for arbitrary noncompact objects.
\paragraph{Editorial transcription note.}
Original author language, mathematics and ordering are preserved. Publisher front matter is replaced by this independent article wrapper. Bibliography is recovered from matching cached publisher HTML, including its TeX formulas. No new proof or mathematical repair is supplied. Surrounding original results are retained for conservative context and are not extra selected corpus claims.
\clearpage
\input{author-body.tex}
\begin{thebibliography}{99}
\bibitem{Balmer05a} P. Balmer - “The spectrum of prime ideals in tensor triangulated categories” , J. reine angew. Math. 588 (2005), p. 149-168
\bibitem{Balmer18a} P. Balmer - “On the surjectivity of the map of spectra associated to a tensor-triangulated functor” , Bull. London Math. Soc. 50 (2018) no. 3, p. 487-495
\bibitem{Balmer19} P. Balmer - “A guide to tensor-triangular classification” , in Handbook of homotopy theory (H. Miller, ed.), Chapman and Hall/CRC, 2019, Available on the author’s web page
\bibitem{Balmer20a} P. Balmer - “Nilpotence theorems via homological residue fields” , Tunis. J. Math. 2 (2020) no. 2, p. 359-378
\bibitem{BalmerCameron20pp} P. Balmer \& J. Cameron - “Computing homological residue fields in algebra and topology” , 2020
\bibitem{BalmerDellAmbrogioSanders16} P. Balmer, I. Dell’Ambrogio \& B. Sanders - “Grothendieck-Neeman duality and the Wirthmüller isomorphism” , Compositio Math. 152 (2016) no. 8, p. 1740-1776
\bibitem{BalmerFavi11} P. Balmer \& G. Favi - “Generalized tensor idempotents and the telescope conjecture” , Proc. London Math. Soc. (3) 102 (2011) no. 6, p. 1161-1185
\bibitem{BensonIyengarKrause08} D. J. Benson, S. B. Iyengar \& H. Krause - “Local cohomology and support for triangulated categories” , Ann. Sci. École Norm. Sup. (4) 41 (2008) no. 4, p. 573-619
\bibitem{BensonIyengarKrause11a} D. J. Benson, S. B. Iyengar \& H. Krause - “Stratifying modular representations of finite groups” , Ann. of Math. (2) 174 (2011) no. 3, p. 1643-1684
\bibitem{BensonIyengarKrause11c} D. J. Benson, S. B. Iyengar \& H. Krause - “Stratifying triangulated categories” , J. Topology 4 (2011) no. 3, p. 641-666
\bibitem{BensonIyengarKrause12b} D. J. Benson, S. B. Iyengar \& H. Krause - “Colocalizing subcategories and cosupport” , J. reine angew. Math. 673 (2012), p. 161-207
\bibitem{BensonIyengarKrause12a} D. J. Benson, S. B. Iyengar \& H. Krause - Representations of finite groups: local cohomology and support , Oberwolfach Seminars , vol. 43 , Birkhäuser/Springer, Basel, 2012
\bibitem{BensonIyengarKrause13} D. J. Benson, S. B. Iyengar \& H. Krause - “Module categories for group algebras over commutative rings” , J. K-Theory 11 (2013) no. 2, p. 297-329, With an appendix by Greg Stevenson
\bibitem{BalmerKrauseStevenson19} P. Balmer, H. Krause \& G. Stevenson - “Tensor-triangular fields: ruminations” , Selecta Math. (N.S.) 25 (2019) no. 1, article ID 13, 36 pages
\bibitem{BalmerKrauseStevenson20} P. Balmer, H. Krause \& G. Stevenson - “The frame of smashing tensor-ideals” , Math. Proc. Cambridge Philos. Soc. 168 (2020) no. 2, p. 323-343
\bibitem{DwyerPalmieri08} W. G. Dwyer \& J. H. Palmieri - “The Bousfield lattice for truncated polynomial algebras” , Homology Homotopy Appl. 10 (2008) no. 1, p. 413-436
\bibitem{HoveyPalmieriStrickland97} M. Hovey, J. H. Palmieri \& N. P. Strickland - Axiomatic stable homotopy theory , Mem. Amer. Math. Soc. , vol. 128, no. 610 , American Mathematical Society, Providence, RI, 1997
\bibitem{HoveyStrickland99} M. Hovey \& N. P. Strickland - Morava $K$-theories and localisation , Mem. Amer. Math. Soc. , vol. 139, no. 666 , American Mathematical Society, Providence, RI, 1999
\bibitem{Krause00} H. Krause - “Smashing subcategories and the telescope conjecture—an algebraic approach” , Invent. Math. 139 (2000) no. 1, p. 99-133
\bibitem{LurieHA} J. Lurie - “Higher algebra” (2017), Online at http://www.math.ias.edu/\textasciitilde{}lurie/
\bibitem{Neeman96} A. Neeman - “The Grothendieck duality theorem via Bousfield’s techniques and Brown representability” , J. Amer. Math. Soc. 9 (1996) no. 1, p. 205-236
\bibitem{Neeman00} A. Neeman - “Oddball Bousfield classes” , Topology 39 (2000) no. 5, p. 931-935
\bibitem{Neeman01} A. Neeman - Triangulated categories , Annals of Math. Studies , vol. 148 , Princeton University Press, Princeton, NJ, 2001
\bibitem{Stevenson13} G. Stevenson - “Support theory via actions of tensor triangulated categories” , J. reine angew. Math. 681 (2013), p. 219-254
\end{thebibliography}
\end{document}
